\documentclass[10pt,reqno]{amsart}
\usepackage[T1]{fontenc}
\usepackage{lmodern,microtype}
\usepackage{amsmath,amssymb,mathtools}
\usepackage[colorlinks=true,linkcolor=blue,citecolor=blue,urlcolor=blue]{hyperref}
\newtheorem{theorem}{Theorem}[section]
\newtheorem{lemma}[theorem]{Lemma}
\newtheorem{corollary}[theorem]{Corollary}
\newtheorem{proposition}[theorem]{Proposition}
\theoremstyle{definition}
\newtheorem{definition}[theorem]{Definition}
\newtheorem{conjecture}[theorem]{Conjecture}
\newtheorem{hypothesis}[theorem]{Hypothesis}
\theoremstyle{remark}
\newtheorem{remark}[theorem]{Remark}
\newcommand{\A}{\mathcal A}
\newcommand{\cP}{\mathcal P}
\newcommand{\cR}{\mathscr R}
\newcommand{\cF}{\mathfrak F}
\newcommand{\cG}{\mathfrak G}
\newcommand{\E}{\mathbb E}
\newcommand{\Prob}{\mathbb P}
\newcommand{\Z}{\mathbb Z}
\newcommand{\ind}{\mathbf 1}
\newcommand{\lcm}{\operatorname{lcm}}
\newcommand{\Bin}{\operatorname{Bin}}
\newcommand{\Bern}{\operatorname{Bern}}
\newcommand{\e}{\mathrm e}
\newcommand{\Qseq}{Q_{\rm seq}}
\newcommand{\lab}[1]{\textup{[\textsc{#1}]}}
\usepackage{mathrsfs}
% TODO(submission): author metadata; decide the citation form for the
% internal repository documents.
\title[Sieve limits for the Erd\H{o}s--Straus exceptional set]{Sieve limits for the Erd\H{o}s--Straus exceptional set:\\ why $3/4$ is sharp for congruence sieves}
\author{Anonymous}
\date{Research note; internally reviewed results only, not externally refereed}
\begin{document}
\begin{abstract}
Let $E(N)$ count the $n\leq N$ for which $4/n$ is not a sum of three
positive unit fractions.  The best internally proved bound in this
programme is $E(N)\ll N\exp\{-c(\log N)^{3/4}\}$.  We show that the
exponent $3/4$ is a hard ceiling for a precisely defined class of
congruence sieves.  The class consists of nonnegative CRT majorants that
are at least one on the whole avoider set of a family of forced
congruence classes, whose final bound has the form
$N\cdot\E\nu+\sum_i|a_i|$, and whose family moduli are at most
$N^{O(1)}$.  For families in which every modulus has a dominant prime
factor, and, for the classes $\cR(M)$, for families of gapped moduli, every
such sieve saves at most $O((\log N)^{3/4})$.  No power of $\log\log N$
can be gained.  The majorant of the $3/4$ note lies in this class, and the
large-sieve argument of the $2/3$-loglog note is capped in the same way at
its own exponent.  We list exactly what the theorems exclude.  For
Selberg-type majorants we prove a sieve limit with no structural
hypothesis, in terms of a noise-stability quantity, and a two-prime cap
with twin moduli that is conditional on an equidistribution hypothesis.
Every statement carries its status label.  All proofs were checked
internally only; nothing here has been externally refereed.  This is not a
proof of the Erd\H{o}s--Straus conjecture.
\end{abstract}
\maketitle

\section{Introduction and status}\label{sec:intro}

Call $n\geq1$ \emph{representable} if $4/n=1/x_1+1/x_2+1/x_3$ with
positive integers $x_i$, and let $E(N)$ count the non-representable
$n\leq N$.  The Erd\H{o}s--Straus conjecture asserts that every $n\geq2$ is
representable.  Vaughan~\cite{Vaughan} proved
$E(N)\ll N\exp\{-c(\log N)^{2/3}\}$ with Montgomery's large
sieve~\cite{Montgomery}.  Two companion notes of this programme
improve this.  The first~\cite{LL} gains a factor
$(\log\log N)^{1/3}$ in the exponent.  The second~\cite{TQ} proves
\[
 E(N)\ll N\exp\{-c(\log N)^{3/4}\},
\]
with label \emph{internally proved; not externally refereed}.  The $3/4$
note itself explains $3/4$ heuristically: the forced classes have
cubic logarithmic mass, and a mass-driven optimisation gives
$\theta=B/(B+1)$ with $B=3$ \cite[\S\ref*{sec:intro}]{TQ}.
% TODO: check the exact section of the 3/4 note for the ceiling heuristic
% (its Section "The heuristic mass-driven three-quarter ceiling").

This note turns that heuristic into a theorem for a precise class of
sieves, and records exactly where the theorem stops.

\subsection*{Status labels}
Every statement carries one of the following labels.
\begin{itemize}
\item \lab{proved}: proved here, with the proof given or precisely cited
  to an internal document; checked by internal hostile review only.
\item \lab{proved mod X}: proved assuming the published external result X,
  which was checked against the source but not re-proved.
\item \lab{conditional}: proved assuming a named open hypothesis.
\item \lab{evidence}: numerical or model evidence; never a proof.
\item \lab{conjecture} / \lab{open}: not proved.
\end{itemize}
Each \lab{proved} result below appears in one of four internal working
documents, \cite{ET}, \cite{EB}, \cite{TW}, \cite{TW2}.  Each was
hostile-reviewed; the reviews are \cite{ETrev,EBrev,TWrev,TW2rev}.  No
prior source for the sieve-limit theorem (Theorem~\ref{thm:slice}) was
found; it may be folklore in spirit.  Novelty is not claimed.

\subsection*{Results}
% TODO: fill in after the body is written.

\section{Forced classes}\label{sec:forced}

A residue class $b\pmod G$ is \emph{forced} if every positive integer
$n\equiv b\pmod G$ is representable.  We use three families of forced
classes.  For an integer $D\geq1$ put
$g(D)=\prod_pp^{\lceil v_p(D)/2\rceil}$, so that $g(D)\mid D\mid g(D)^2$
and, for every integer $x$, $D\mid x^2\iff g(D)\mid x$.

\begin{lemma}[multiplier identity; {\cite[Lemma~2.1]{LL}}; \lab{proved}]\label{lem:identity}
Let $M\equiv3\pmod4$, $A=(M+1)/4$, and $A=uvw$ with positive integers
$u,v,w$.  If $n\geq1$ and $nv\equiv-u\pmod M$, then $s=(nv+u)/M$ is a
positive integer and $4/n=1/(suw)+1/(nsvw)+1/(nuvw)$.
\end{lemma}
\begin{proof}
On the common denominator $nsuvw$ the numerators sum to
$nv+u+s=s(M+1)=4suvw$.
\end{proof}

\begin{lemma}[the classes $\cR(M)$; {\cite[Lemma~18.1]{Notes}}; \lab{proved}]\label{lem:RM}
For $M\equiv3\pmod4$ let $\cR(M)$ be the set of classes $-uv^{-1}\pmod M$
over all factorisations $A=uvw$.  Then
\[
 \cR(M)=\{-4D\pmod M:\ D\mid A^2\},\qquad
 |\cR(M)|\leq\tau(A^2),
\]
and every class of $\cR(M)$ is forced.
\end{lemma}
\begin{proof}
Since $4A\equiv1\pmod M$, a factorisation $A=uvw$ gives
$-uv^{-1}\equiv-u^2wA^{-1}\equiv-4u^2w\pmod M$ with $u^2w\mid A^2$.
Conversely, if $A=\prod p^{e_p}$ and $D=\prod p^{d_p}$ with $d_p\leq2e_p$, put
$v_p(u)=\lfloor d_p/2\rfloor$, $v_p(w)=d_p-2\lfloor d_p/2\rfloor$ and
$v_p(v)=e_p-v_p(u)-v_p(w)\geq0$; then $uvw=A$ and $u^2w=D$.  Forcedness is
Lemma~\ref{lem:identity}.
\end{proof}

\begin{lemma}[$(a,D)$-classes; {\cite[Lemma~3.2]{ET}}; \lab{proved}]\label{lem:aD}
For all integers $a,D\geq1$ the class $n\equiv-(4D+a)\pmod{4a\,g(D)}$ is
forced.
\end{lemma}
\begin{proof}
Write $n+4D+a=4a\,g(D)\,j$ with $j\geq1$.  Then $M:=(n+4D)/a=4g(D)j-1$ is
$\equiv3\pmod4$ and $(M+1)/4=g(D)j$, so $D\mid g(D)^2\mid((M+1)/4)^2$.
Since $n\equiv-4D\pmod M$, $n$ lies in a class of $\cR(M)$, and
Lemma~\ref{lem:RM} applies.
\end{proof}
These classes encode exactly the Case-B witnesses of the classical
parametrisation \cite[Thm~3.1(B)]{Notes} (see \cite[Lemma~3.2]{ET}).

\begin{lemma}[Case-A classes; {\cite[\S3]{ET}}; \lab{proved}]\label{lem:caseA}
For $d\geq1$ and $m\mid4d+1$ (so $m$ is odd and $(m,4g(d))=1$), the class
$n\equiv-m^{-1}\pmod{4g(d)}$ is forced.
\end{lemma}
\begin{proof}
Put $z_0=(nm+1)/4$, an integer.  Then $g(d)\mid z_0$, so $d\mid z_0^2$;
and $m\mid 4d+1$ gives $m\mid 4d+4z_0$, so $m\mid d+z_0$.  With
$x=(z_0+d)/m$ and $y=(z_0+z_0^2/d)/m=z_0(d+z_0)/(dm)$, both integers since
$(m,d)=1$, one has $1/x+1/y=m/z_0$, hence
$1/x+1/y+1/(nz_0)=(nm+1)/(nz_0)=4/n$.
\end{proof}

We write \emph{Case-B classes} for the classes of $\cR(M)$ and the
$(a,D)$-classes, and \emph{Case-A classes} for those of
Lemma~\ref{lem:caseA}.  Two facts are used repeatedly.

\begin{lemma}[$n=1$ is never covered; \lab{proved}]\label{lem:one}
No forced class contains $n=1$.
\end{lemma}
\begin{proof}
$4/1>3\geq1/x_1+1/x_2+1/x_3$.
\end{proof}

\begin{lemma}[Mordell's obstruction; classical; {\cite[Lemma~1.1]{TW}}; \lab{proved}]\label{lem:mordell}
Let $M\equiv3\pmod4$, $A=(M+1)/4$ and $D\mid A^2$.  Then $(D,M)=1$ and
the Jacobi symbol satisfies $\bigl(\frac{-4D}{M}\bigr)=-1$.  Consequently,
if $n$ is a nonzero quadratic residue modulo every prime factor of $M$,
then $n$ lies in no class of $\cR(M)$.
\end{lemma}
\begin{proof}
$(A,M)=(A,4A-1)=1$.  Write $D=sr^2$ with $s$ squarefree; then $D\mid A^2$
gives $sr\mid A$, so $M\equiv-1\pmod{4s}$.  Now
$(\frac{-4D}M)=(\frac{-1}M)(\frac sM)=-(\frac sM)$.  If $s$ is odd,
reciprocity and $(\frac Ms)=(\frac{-1}s)$ give $(\frac sM)=1$, since
$(M-1)/2$ is odd.  If $s=2s'$, then $8\mid M+1$, so $(\frac2M)=1$, and
$(\frac{s'}M)=1$ as before.  For the consequence,
$(\frac nM)=\prod_{p^e\| M}(\frac np)^e=1$.
\end{proof}

\end{document}
